\section{Path rectangles and the universal cover}\label{sec:global}

We now turn the holomorphic box theorem into continuation along arbitrary
continuous paths. The important point is that a transported edge retains
a holomorphic parametrization on an entire polydisk. Its image need not
remain small or lie in a single original chart.

\begin{proposition}[Path rectangles]\label{prop:rectangles}
The following assertions hold.
\begin{enumerate}
\item If \(\alpha,\beta:[0,1]\to X\) locally follow \(E_1\)- and
\(E_2\)-plaques, respectively, and \(\alpha(0)=\beta(0)\), there is
a continuous box \(R:[0,1]^2\to X\) with
\(R(s,0)=\alpha(s)\) and \(R(0,t)=\beta(t)\).
\item Every continuous path \(\gamma:[0,1]\to X\) is the diagonal
of a continuous box: \(R(s,s)=\gamma(s)\).
\end{enumerate}
\end{proposition}

\begin{proof}
For the first assertion, subdivide the horizontal interval into pieces
\(I_i=[u_{i-1},u_i]\) and the vertical interval into pieces
\(J_j=[v_{j-1},v_j]\) on which the respective paths lie in single
plaques. Each horizontal piece has a representation
\[
 \alpha(s)=A_i(a_i(s)),\qquad s\in I_i,
\]
where \(a_i:I_i\to\D^{r_1}\) is continuous and
\(A_i:\D^{r_1}\to X\) is a holomorphic plaque parametrization
tangent to \(E_1\). Similarly write
\(\beta(t)=B_j(b_j(t))\), with
\(b_j:J_j\to\D^{r_2}\) continuous and \(B_j\) holomorphic and
tangent to \(E_2\). Here \(r_i=\rank E_i\).

Fill the tiles \(I_i\times J_j\) from the bottom and left. Inductively,
the known edges of a tile have the forms
\(A(a_i(s))\) and \(B(b_j(t))\), where \(A\) and \(B\) are
holomorphic on the \emph{entire} polydisks \(\D^{r_1}\) and
\(\D^{r_2}\), tangent to the respective bundles. They agree at their
common corner. Apply Proposition~\ref{prop:polydisk-axes} with origins
\(a_i(u_{i-1})\) and \(b_j(v_{j-1})\), obtaining a holomorphic box
\(G:\D^{r_1}\times\D^{r_2}\to X\). Its evaluation
\[
 R(s,t)=G(a_i(s),b_j(t))
\]
fills the tile. The opposite horizontal and vertical edges have the
holomorphic representatives
\[
 a\longmapsto G(a,b_j(v_j)),\qquad
 b\longmapsto G(a_i(u_i),b),
\]
again on the whole original polydisks. This proves the induction without
any shrinking of the edge domains.

For the second assertion, choose a subdivision
\(0=u_0<\cdots<u_N=1\) such that each
\(\gamma(I_i)\), \(I_i=[u_{i-1},u_i]\), lies in a product chart
with inverse \(\theta_i:\D^{r_1}\times\D^{r_2}\to X\).
Write \(\gamma(s)=\theta_i(a_i(s),b_i(s))\) on that interval.
Define the diagonal tile by
\[
 R(s,t)=\theta_i(a_i(s),b_i(t)),\qquad (s,t)\in I_i^2.
\]
These tiles agree at consecutive diagonal vertices, and their edges have
the full-polydisk representations used above.

Fill the remaining tiles \(I_i\times I_j\) in successive bands away
from the diagonal, that is, in increasing order of \(|i-j|\).
If \(i<j\), use the already filled bottom and right
edges; if \(i>j\), use the top and left edges. The same polydisk box
construction works from any corner, since the origins in
Proposition~\ref{prop:polydisk-axes} are arbitrary. The two input edges
have matching corner values: for \(|i-j|=1\) this is the prescribed
diagonal value, and for \(|i-j|>1\) it follows from the intervening
tile closer to the diagonal. More explicitly, when \(i<j-1\), the
bottom-right corner is shared with the previously filled tile
\(I_{i+1}\times I_{j-1}\). The case \(i>j+1\) is symmetric.
Tiles filled at the same stage share no edge, and all their shared
vertices already lie on prescribed input edges. Thus every tile is
consistent with the previously assigned values, and its output edges
retain their full-polydisk representatives.

In either construction, the finite collection of tiles glues
continuously. The local plaque condition holds across seams for the
same reason as in the proof of Proposition~\ref{prop:polydisk-axes}:
in a product chart near a seam point, the unwanted coordinate is
constant on consecutive slice pieces and agrees at their common
endpoints. The resulting map is a continuous box.
\end{proof}

Figure~\ref{fig:grid} illustrates both filling orders. The transported
edges retain their full-polydisk representatives; this, rather than a
bound on their image sizes, is what allows each new tile to be filled.
\begin{figure}[tbp]
\centering
\setlength{\unitlength}{1pt}
\begin{picture}(430,175)
  \put(102,163){\makebox(0,0){\small (a) prescribed axes}}
  \put(322,163){\makebox(0,0){\small (b) prescribed diagonal}}
  {\color{black!12}
   \put(55,45){\rule{30pt}{30pt}}
   \put(275,45){\rule{30pt}{30pt}}
   \put(305,75){\rule{30pt}{30pt}}
   \put(335,105){\rule{30pt}{30pt}}}\ignorespaces
  {\color{blue!12}
   \put(85,45){\rule{30pt}{30pt}}
   \put(305,45){\rule{30pt}{30pt}}}\ignorespaces
  \thinlines
  \multiput(55,45)(30,0){4}{\line(0,1){90}}
  \multiput(55,45)(0,30){4}{\line(1,0){90}}
  \multiput(275,45)(30,0){4}{\line(0,1){90}}
  \multiput(275,45)(0,30){4}{\line(1,0){90}}
  {\color{black!55}\put(275,45){\line(1,1){90}}}\ignorespaces
  \thicklines
  {\color{blue!65!black}
   \put(55,45){\line(1,0){90}}
   \put(85,45){\vector(1,0){30}}
   \put(305,75){\vector(1,0){30}}}\ignorespaces
  {\color{red!65!black}
   \put(55,45){\line(0,1){90}}
   \put(85,45){\vector(0,1){30}}
   \put(305,75){\vector(0,-1){30}}}\ignorespaces
  \put(85,45){\circle*{3}}
  \put(305,75){\circle*{3}}
  \put(100,60){\makebox(0,0){\small $R$}}
  \put(320,60){\makebox(0,0){\small $R$}}
  \put(100,30){\makebox(0,0){\small $\alpha(s)$}}
  \put(47,90){\makebox(0,0)[r]{\small $\beta(t)$}}
  \put(155,45){\makebox(0,0){\small $s$}}
  \put(55,144){\makebox(0,0){\small $t$}}
  \put(375,45){\makebox(0,0){\small $s$}}
  \put(275,144){\makebox(0,0){\small $t$}}
  \put(320,29){\makebox(0,0){\small $R(s,s)=\gamma(s)$}}
  {\color{black!12}\put(105,5){\rule{9pt}{9pt}}}\ignorespaces
  \put(120,9){\makebox(0,0)[l]{\small completed tiles}}
  {\color{blue!12}\put(255,5){\rule{9pt}{9pt}}}\ignorespaces
  \put(270,9){\makebox(0,0)[l]{\small next tile}}
\end{picture}
\caption{The two grid constructions in the parameter square. In (a),
the bottom and left axes are prescribed, and tiles are filled in
increasing order of the sum of their indices. In (b), the diagonal tiles
are prescribed first, and the remaining tiles are filled in successive
bands away from the diagonal. In each highlighted tile, the blue
horizontal edge and red vertical edge are already known and meet at
the marked corner. The box with these axes supplies the other two edges.
The reversed red arrow in (b) indicates that the same construction can
start at any corner. No target-space size estimate is depicted.}
\label{fig:grid}
\end{figure}


Fix a basepoint \(o\in X\). For \(i=1,2\), let \(L_i\) be the
leaf through \(o\), endowed with its intrinsic plaque topology, not
the subspace topology inherited from \(X\).

\begin{lemma}[Intrinsic leaves]\label{lem:leaves}
Each \(L_i\) is a connected, Hausdorff, second-countable complex
manifold of dimension \(r_i\), and its inclusion into \(X\) is an
injective holomorphic immersion. It admits a connected simply connected
cover \(\widetilde L_i\), also a second-countable complex manifold.
A path locally following \(E_i\)-plaques is continuous in the intrinsic
topology of its leaf.
\end{lemma}

\begin{proof}
The leaf consists of points reachable from \(o\) by finite chains of
\(E_i\)-plaques. Plaque coordinates give its complex atlas: near a
point of intersection, two plaques through that point coincide, and
their coordinate change is holomorphic. The inclusion into \(X\) is
an injective continuous immersion. In particular, distinct points of
the leaf are separated by inverse images of disjoint ambient open sets,
so the leaf is Hausdorff. The definition by plaque chains gives path
connectedness and the asserted continuity of plaque paths.

For second countability, take a countable ambient product atlas. Start
with a plaque through \(o\). For each plaque already reached and each
ambient chart, their intersection has at most countably many connected
components: it is an open subset of a second-countable plaque. Each
component lies in a single plaque of the new chart, since the transverse
coordinate is locally constant. Adjoin these plaques and repeat.
There are only countably many plaques after all finite stages, and
every plaque path from \(o\) is covered by a finite chain of this
kind. Their countable coordinate bases form a countable basis for the
leaf.

The usual covering-space construction applies because a manifold is
locally path connected and semilocally simply connected. Its universal
cover is second countable as well. One way to see the latter detail is
to refine the countable atlas to simply connected coordinate disks.
There are countably many components of their pairwise overlaps. A
based loop is described, up to homotopy, by a finite sequence of these
charts and overlap components; fixing paths within the charts and
components supplies representatives for all such descriptions. Hence
the fundamental group is countable. The lifted countable chart basis
has only countably many sheets and is therefore a countable basis for
the universal cover. Lifting the holomorphic charts gives its complex
structure.
\end{proof}

A product chart at \(o\) supplies mutually inverse germs
\[
 H_o:(L_1\times L_2,(o,o))\longrightarrow(X,o),
 \qquad
 K_o:(X,o)\longrightarrow(L_1\times L_2,(o,o)).
\]
The first combines coordinates on the two plaques through \(o\); the
second projects onto them. Separated chart transitions show that the
germs do not depend on the chosen product chart. The next result gives
analytic continuation in both directions.

\begin{proposition}[Transport of inverse germs]\label{prop:germ-transport}
Let \(R:[0,1]^2\to X\) be a continuous box with \(R(0,0)=o\), and set
\[
 \alpha(r)=R(r,0)\in L_1,
 \qquad \beta(r)=R(0,r)\in L_2,
 \qquad \gamma(r)=R(r,r)\in X.
\]
There are mutually inverse local biholomorphism germs \(H_r,K_r\)
at \((\alpha(r),\beta(r))\) and \(\gamma(r)\), respectively,
which give analytic continuation of \(H_o\) along \((\alpha,\beta)\)
and of \(K_o\) along \(\gamma\).
\end{proposition}

\begin{proof}
Transport a first-coordinate germ from the \(L_1\)-plaque at
\(\alpha(r)\) along the vertical path
\(u\mapsto R(r,ru)\), \(0\leq u\leq1\). To define this transport,
subdivide the path into finitely many subsegments lying in product
charts, and choose a local separated transition at each junction.
Within a subsegment the first coordinate is constant; at a junction it
changes by the first-coordinate local biholomorphism of that transition.
Their composition transports a variable first coordinate to the
endpoint. Refining the subdivision or changing intermediate charts
does not change the germ: separated transitions satisfy the cocycle
identity. On a common vertical subsegment, the first-coordinate transition
germ at the fixed first coordinate is locally constant, since local
separated transitions are independent of the second coordinate. It is
therefore constant along the connected subsegment. Common refinements identify any two
constructions. Similarly, transport a second-coordinate germ from the
\(L_2\)-plaque at \(\beta(r)\) along
\(u\mapsto R(ru,r)\). Combining the two transported coordinates in
one product chart at \(\gamma(r)\) gives a local biholomorphism germ
\(H_r\); let \(K_r\) be its inverse.

We must show genuine analytic continuation, rather than merely
continuous dependence of a collection of germs. Fix \(r_0\). Choose
the finite chart chains for the two reference transport paths, a fixed
small overlap neighborhood at each junction, and a common final chart
\(z\) around \(\gamma(r_0)\). By joint continuity of \(R\), the
same compact subsegments lie in the same charts, and the same junctions
lie in the same overlap neighborhoods, for all \(r\) sufficiently
near \(r_0\). Also \(\alpha(r)\) and \(\beta(r)\) remain in
fixed intrinsic leaf plaques by Lemma~\ref{lem:leaves}.

Use coordinates \(a\) and \(b\) on those two starting plaques.
Every first-coordinate transition in the vertical chain is now one
fixed holomorphic map, independent of the second coordinate. After
shrinking the initial coordinate neighborhood, their finite composition
is a single local biholomorphism \(A\) defined near
\(a(\alpha(r_0))\). It describes the transport for every nearby
\(r\), not just for \(r_0\). The horizontal chain similarly gives
one fixed second-coordinate local biholomorphism \(B\). Consequently
one holomorphic map
\begin{equation}\label{eq:fixed-transport}
 (x,y)\longmapsto
 z^{-1}\bigl(A(a(x)),B(b(y))\bigr)
\end{equation}
represents \(H_r\) at \((\alpha(r),\beta(r))\) for every nearby
\(r\), and sends that point to \(\gamma(r)\). After one further
shrinking, its single holomorphic inverse represents all the nearby
\(K_r\). At \(r=0\), both transport paths are constant, so these
germs are precisely \(H_o\) and \(K_o\). This is the required
analytic continuation.
\end{proof}

\begin{proof}[Proof of Theorem~\ref{thm:main}]
Every path from \((o,o)\) in \(L_1\times L_2\) consists of two
intrinsically continuous leaf paths. Proposition~\ref{prop:rectangles}
provides a box with these as its axes, and
Proposition~\ref{prop:germ-transport} continues \(H_o\) along the
given path. Conversely, every path from \(o\) in \(X\) is the
diagonal of a box, so the same two Propositions continue \(K_o\)
along every such path.

Let \(\pi:\widetilde X\to X\) and
\(q_i:Y_i=\widetilde L_i\to L_i\) be the based connected simply
connected covers, and put \(q=q_1\times q_2\). The product
\(Y_1\times Y_2\) is the universal cover of \(L_1\times L_2\).
The monodromy theorem, applied on these simply connected source covers,
gives holomorphic maps
\[
 H:Y_1\times Y_2\longrightarrow X,
 \qquad K:\widetilde X\longrightarrow L_1\times L_2
\]
with the prescribed initial germs. Here the usual monodromy argument
applies to manifold-valued holomorphic germs: uniqueness follows from
the identity theorem, and finite chains of local representatives are
stable under small path variations, so homotopic paths with fixed
endpoints produce the same terminal germ.

Since both source manifolds are simply connected, these maps lift to
holomorphic maps
\[
 \widetilde H:Y_1\times Y_2\longrightarrow\widetilde X,
 \qquad
 \widetilde K:\widetilde X\longrightarrow Y_1\times Y_2.
\]
Choose the lifts that send the based points to one another. Their germs
there are mutually inverse, because \(H_o\) and \(K_o\) were inverse.
Thus \(\widetilde H\circ\widetilde K\) and
\(\widetilde K\circ\widetilde H\) agree with the respective identity
maps on neighborhoods of the basepoints. The identity theorem on the
connected source manifolds makes both compositions identities
everywhere. In particular, \(\Phi=\widetilde K\) is a biholomorphism
\(\widetilde X\simeq Y_1\times Y_2\).

It remains to retain the specified splitting. Near the basepoint, the
initial product chart gives
\[
 \dd\Phi(\pi^*E_1)\subset\pr_1^*T_{Y_1},\qquad
 \dd\Phi(\pi^*E_2)\subset\pr_2^*T_{Y_2}.
\]
Equivalently, the two wrong-factor components of \(\dd\Phi\) vanish
there. They are holomorphic bundle maps, so they vanish on all of the
connected manifold \(\widetilde X\). Finally, \(\dd\Phi\) is
invertible and \(\dim_\C Y_i=r_i\) by Lemma~\ref{lem:leaves}; hence
these inclusions are equalities. The factors are connected and simply
connected by construction, proving the asserted compatible splitting.
\end{proof}
